% Personal project documentation, Section 19 of the build specification.
%
% Written for my own archive: the record I would want if I picked this up again
% in two years having forgotten why anything is the way it is. Detail is the
% point, not concision.
%
% Ground rule 1 applies here as everywhere: no double hyphen in prose.

\documentclass[11pt,a4paper]{report}

\usepackage[T1]{fontenc}
\usepackage[utf8]{inputenc}
\usepackage[english]{babel}
\usepackage{lmodern}
\usepackage{microtype}
\usepackage{geometry}
\geometry{margin=2.5cm}

\usepackage{amsmath}
\usepackage{amssymb}
\usepackage{booktabs}
\usepackage{longtable}
\usepackage{listings}
\usepackage{xcolor}
\usepackage{caption}
\usepackage{enumitem}
\usepackage[hidelinks]{hyperref}

\captionsetup{font=small,labelfont=bf}
\setlength{\parskip}{0.4em}
\setlength{\parindent}{0pt}

\definecolor{codebg}{HTML}{F5F5F2}
\lstset{
  basicstyle=\ttfamily\footnotesize,
  backgroundcolor=\color{codebg},
  frame=none,
  breaklines=true,
  showstringspaces=false
}

\title{\textbf{Parallel Numerical Library}\\[0.4em]
  \large Complete project documentation, for my own reference}
\author{Olajide Badejo}
\date{August 2026}

\begin{document}

\maketitle
\tableofcontents

% =====================================================================
\chapter{Project overview}

\section{What this is}

A C++23 numerical methods library with a pluggable parallel execution backend
layer. Twelve iterative linear solvers are implemented exactly once, against a
single problem interface, and run unchanged over seven execution backends:
serial, OpenMP, POSIX threads, a \texttt{std::jthread} pool, MPI, hybrid MPI with
OpenMP threads, and CUDA.

Around that core sit a numerics module (root finding, quadrature, ODE
integration, dense LU and QR), two problem families (the 2D Poisson five point
stencil and seeded dense systems), a benchmark harness, and three reports.

\section{Why I built it this way}

The project merges two earlier ideas, a numerical methods library and a parallel
scientific computing study, and the merge is only worth doing if it produces
something neither would alone. The thing it produces is this: \textbf{one
numerical library, many execution backends}, so that a benchmark can isolate
what each programming model costs on identical numerical work.

The usual comparison between parallel programming models is compromised before
it starts, because each model gets its own implementation and the study ends up
measuring the care that went into each one. Here a solver does not know which
backend it is on and a backend does not know which solver it is running, so the
only variable between two measurements is the execution model.

That only works if "identical numerical work" is literally true, which is why the
central engineering commitment of the project is bit identical results across
backends, verified by test rather than asserted.

\section{Goals}

\begin{enumerate}[leftmargin=*]
  \item A numerical core where every result carries its own evidence: value,
        error estimate, iteration count, evaluation count, converged flag and an
        explicit stop reason. No caller can consume an unconverged answer by
        accident.
  \item A solver zoo presented as one family tree, with each method derived from
        the common splitting framework and each convergence claim checked against
        closed form theory rather than asserted.
  \item Execution backends behind one interface, with identical numerics as a
        tested invariant.
  \item A backend cost study: the same solver on the same seeded problem across
        every backend, with scaling curves and the heterogeneous core knee.
  \item A CPU versus GPU comparison designed so the hardware does not decide the
        argument.
  \item Professional finish: tests, docs, CI, three PDFs, all reproducible from a
        clean clone with \texttt{make all}.
\end{enumerate}

\section{Scope, and what is deliberately outside it}

Inside: iterative linear solvers on two problem families, seven execution
models, bandwidth bound performance analysis, and convergence theory for the
model problem.

Outside, and deliberately so:

\begin{itemize}[leftmargin=*]
  \item \textbf{Sparse matrix formats.} The stencil is matrix free and the dense
        problem is dense. A general sparse format would add a large amount of
        code that the questions being asked do not need.
  \item \textbf{Preconditioned conjugate gradient.} The pieces exist and it is
        listed under further work, but it was not required and the project was
        already large.
  \item \textbf{Multigrid.} Same reasoning. It would change the iteration count
        story entirely and deserves its own project.
  \item \textbf{Compute bound kernels.} Every device conclusion here is about
        bandwidth bound stencil work, and I say so wherever a conclusion is
        stated.
  \item \textbf{Power, cost, and programmer time.} Not measured, not claimed.
\end{itemize}

% =====================================================================
\chapter{Specifications}

\section{Target machine}

\begin{center}
\begin{tabular}{ll}
\toprule
Component & Specification \\
\midrule
CPU & Intel Core i7-14700K, 8 performance cores plus 12 efficiency cores, 28 threads \\
RAM & 31.7 GB host, 12 GB budgeted to the WSL2 guest \\
GPU & NVIDIA GeForce RTX 5070, 12227 MiB, compute capability 12.0, 48 SMs \\
Driver & 610.62 \\
OS & Windows 11 Pro 10.0.26200, everything running inside WSL2 \\
Guest & Ubuntu 26.04 LTS, kernel 6.18.33.2-microsoft-standard-WSL2 \\
\bottomrule
\end{tabular}
\end{center}

\section{Toolchain as actually installed}

Every version below was verified by running the tool, not by reading a package
list. Where the installed version differs from what the specification asked for,
the substitution and its reason are recorded.

\begin{center}
\begin{longtable}{lll}
\toprule
Component & Specification floor & Installed \\
\midrule
\endhead
Host compiler & GCC 16.1 & GCC 16.0.1 20260322 (experimental) \\
C++ standard & C++23 & \texttt{\_\_cplusplus} 202302 \\
CMake & 4.4 & 4.4.0 via pipx \\
MPI & OpenMPI 5.0.x & OpenMPI 5.0.10 \\
OpenMP & 6.0 cited & \texttt{\_OPENMP} 202111, which is 5.2 \\
CUDA & 13.3 & 13.3, V13.3.73 \\
CUDA host compiler & not specified & GCC 14 \\
Python & 3 with matplotlib, pandas & 3.14.4, matplotlib 3.10.7, pandas 2.3.3 \\
LaTeX & TeX Live with latexmk & latexmk 4.87 \\
Build tool & not specified & Ninja 1.13.2 \\
Lint & ruff & 0.15.22 \\
Format & clang-format & 20.1.7 \\
\bottomrule
\end{longtable}
\end{center}

\subsection{Why GCC 16 over the stable GCC 15}

GCC 16.0.1 is an experimental trunk snapshot and GCC 15.2.0 is a stable release,
so the safe choice looked like 15. I tested both against the features the project
actually needs and chose 16 on merit rather than on version number:

\begin{itemize}[leftmargin=*]
  \item GCC 16 provides \texttt{<mdspan>}, which GCC 15 does not. That is
        directly useful for a 2D grid.
  \item GCC 16 reports \texttt{\_OPENMP} 202111, which is OpenMP 5.2. GCC 15
        reports 201511, which is 4.5.
  \item Both compile the full C++23 feature set the project uses:
        \texttt{jthread}, \texttt{barrier}, \texttt{expected}, ranges,
        \texttt{print}.
\end{itemize}

\subsection{Deliberate deviations from the specification}

\paragraph{WSL processor count.} The specification asks for
\texttt{processors=20} in \texttt{.wslconfig}. The file already set
\texttt{processors=28}, exposing every host logical processor, which is a
superset of what a sweep to twenty workers needs. It also carries a dated comment
explaining that a lower count makes the hybrid core comparison impossible to
measure, and documents a machine crash that the memory budget in the same file
exists to prevent. I left it alone.

\paragraph{OpenMP level.} The 6.0 specification is cited as the reference
document, as asked, but GCC implements to 5.2 and the code is restricted to that.
The build asserts \texttt{\_OPENMP >= 201511} so a downgraded toolchain fails
loudly rather than silently producing different code.

\section{Constraints that shaped the design}

\begin{enumerate}[leftmargin=*]
  \item \textbf{No em dashes or en dashes anywhere, in any file type, including
        compiled PDFs.} Enforced by a linter from the first commit. This also
        rules out the LaTeX \verb|--| and \verb|---| ligatures in prose, and
        BibTeX page ranges, both of which typeset as the banned characters.
  \item \textbf{Reproducibility is a deliverable.} A clone plus
        \texttt{make setup \&\& make all} reproduces every number. Seeds
        recorded, nothing hand copied.
  \item \textbf{No number without a run.} Pending values read \texttt{pending}.
        Every result row carries backend, worker count, pinning, seed and commit
        hash.
  \item \textbf{Honest comparison.} The report never headlines a GPU speedup
        without the device normalised framing beside it.
  \item \textbf{12 GB memory ceiling in the guest.} This bounds problem sizes and
        is why build parallelism is capped at six jobs.
\end{enumerate}

% =====================================================================
\chapter{Implementation details}

\section{Core types and diagnostics}

\texttt{include/pnl/core/} holds the vocabulary. \texttt{Real} is
\texttt{double} throughout; there is no mixed precision path, so a single alias
keeps the CUDA boundary unambiguous. \texttt{Index} is \texttt{std::ptrdiff\_t},
deliberately signed: OpenMP canonical loop form accepts unsigned indices only
from 3.0 onward, several compilers still generate worse code for them, and a
backward sweep is expressible without wraparound traps.

\texttt{Diagnostics} is the record attached to every numerical result: error
estimate, iterations, evaluations, converged flag, and a \texttt{StopReason}
enumeration distinguishing convergence from hitting the cap, stagnation,
breakdown and divergence. Hitting an iteration cap is never indistinguishable
from meeting a tolerance.

\texttt{block\_partition(n, parts, k)} lives here rather than inside a backend,
because both the chunk grid and the MPI row decomposition use it. It is remainder
aware: the first \texttt{n mod parts} blocks get one extra element, so sizes
differ by at most one and nothing is dropped.

\section{The backend interface, and the decision it rests on}

\texttt{backend.hpp} declares \texttt{parallel\_for}, \texttt{reduce},
\texttt{barrier}, plus \texttt{local\_rows}, \texttt{exchange\_halo},
\texttt{gather\_rows} and \texttt{run\_ordered} for the distributed case. It is
the minimal common denominator of the models it spans.

\subsection{Chunk level callbacks}

The body of a \texttt{parallel\_for} receives a \texttt{Range} and loops over it
itself, rather than being called once per element. This was the first significant
design decision and it has two consequences I wanted:

\begin{itemize}[leftmargin=*]
  \item The \texttt{std::function} indirection is paid once per chunk, so it is
        proportional to the worker count and not to the problem size. It does not
        contaminate the timings.
  \item The innermost loop stays a plain loop over contiguous memory that the
        compiler can vectorise, so each backend measures its own dispatch cost
        rather than a penalty the abstraction imposed on everything equally.
\end{itemize}

\subsection{Determinism, which turned out to be the keystone}

Floating point addition is not associative, so a reduction's answer depends on
the grouping of its partial sums. Most libraries treat this as unavoidable. I
made it a property of the design instead.

Under \texttt{ReductionMode::Deterministic} the chunk grid is fixed by the
problem size alone, \texttt{min(512, n)} chunks, never by the worker count. Each
chunk's partial goes into its own slot and the slots are summed in ascending
index order. Which worker computed which slot, and when it finished, cannot
affect the result.

The payoff is larger than I expected when I wrote it. Because every shared memory
backend produces bit identical iterates at every worker count, the equivalence
suite can assert \emph{exact equality} rather than a tolerance. That makes it a
far more sensitive instrument: it fails on the first bit that differs. The fused
multiply add discrepancy described later was one bit in the last place and any
reasonable tolerance would have swallowed it.

\section{Problems}

\subsection{Poisson2D}

The five point stencil on the unit square. Stored as an $(n+2)$ by $(n+2)$ row
major grid with a one cell boundary ring holding the homogeneous Dirichlet
values, so there is no boundary branch in the innermost loop and the ring doubles
as the MPI halo. The scaling by $h^2$ makes the matrix have diagonal 4 and off
diagonals $-1$.

\texttt{PoissonTheory} computes the closed form spectral quantities so that no
test re-derives them: $\rho_J = \cos(\pi h)$,
$\rho_{GS} = \rho_J^2$, $\omega^{*} = 2/(1 + \sin(\pi h))$.

Two right hand sides exist and the reason is important enough that it has its own
section under problems below.

\subsection{DenseProblem}

Seeded dense systems in two flavours: strictly diagonally dominant (guarantees
Jacobi and Gauss Seidel converge, but not symmetric so conjugate gradient does
not apply) and symmetric with a dominant positive diagonal (positive definite by
Gershgorin, so admissible for conjugate gradient and for the Ostrowski and Reich
theorem).

The exact solution is chosen first and the right hand side formed from it, so the
unit tests compare against something that is not another solver's output.
Diagonal blocks are LU factorised once in the constructor and reused every sweep,
which is what makes the block methods competitive rather than merely correct.

\section{The solver zoo}

Twelve solvers, eleven of which are one splitting family. Each lives in its own
header, documents its splitting, its convergence order and the exceptions it
throws, and hands a sweep to a shared driver. The shared driver is what
guarantees every method measures its residual the same way, stops on the same
criterion, and pays the same reporting overhead.

The block methods take a block count as a \emph{solver parameter and never the
worker count}, which is what keeps their iterates independent of how many workers
happen to be running. Each problem defines its natural block: a grid line solved
by the Thomas algorithm for the stencil, contiguous index ranges solved by dense
LU for the dense problem.

\section{What each backend does underneath}

\begin{itemize}[leftmargin=*]
  \item \textbf{serial} walks the same chunk grid as the parallel backends
        rather than looping over the whole range, so "serial equals parallel bit
        for bit" is a statement about the parallel backends and not an artefact.
  \item \textbf{openmp} is the only backend that does not partition the range
        itself; it hands the chunk grid to OpenMP's worksharing construct, so
        what is measured is OpenMP's scheduler.
  \item \textbf{pthreads} uses raw POSIX primitives and a persistent pool.
        Workers wait on a \emph{generation counter} rather than a flag, which
        removes the lost wakeup race a boolean would have.
  \item \textbf{jthread} uses \texttt{std::jthread}, \texttt{std::stop\_token}
        and two \texttt{std::barrier} phases in strict alternation. Deliberately
        no work stealing; the schedule cost measurement is the evidence.
  \item \textbf{mpi} uses remainder aware row decomposition and two
        \texttt{MPI\_Sendrecv} calls per halo exchange. Each rank allocates the
        whole grid and computes its own band, which spends memory to keep the
        solver code identical across backends; the communication volume is
        unaffected.
  \item \textbf{hybrid} inherits from the MPI backend rather than
        reimplementing, so any difference the sweep measures between them is the
        threading and nothing else.
  \item \textbf{cuda} is a separate device solve path, not a \texttt{Backend}.
        The state stays in device memory for the whole solve; forcing it behind
        \texttt{parallel\_for} would measure PCIe latency.
\end{itemize}

\section{Instrumentation}

The sweep harness resolves \texttt{benchmarks/sweep\_matrix.yaml}, which declares
nine blocks each carrying a \texttt{why} field saying what question it answers.
The driver is resumable, records provenance on every row, merges atomically
through a temporary file and rename, and writes a session manifest containing
both devices' bandwidth probes and the toolchain versions.

Figures and tables are generated from the summary by
\texttt{scripts/gen\_report\_assets.py}, which is idempotent: the only inputs are
the summary and the manifest, and every output is overwritten in full.

% =====================================================================
\chapter{Chronological account}

\section{Phase 0: environment}

Started by probing rather than assuming. Found the hardware matched the
specification exactly, and that a \texttt{.wslconfig} already existed with a
documented crash history, which I left alone.

Found GCC 16 available in apt and tested it against GCC 15 on the features the
project needs, choosing 16 for \texttt{<mdspan>} and OpenMP 5.2. Installed CMake
4.4 through pipx since the distribution ships 4.2.3.

Probed MPI and found the version macro says 3.1 while MPI 4.0 entry points link.
Probed CUDA and hit the first structural problem: nvcc rejects GCC 16, and cannot
parse GCC 15's headers either. Resolved with a C ABI boundary.

Spent time on a false lead here: the CUDA configure failure looked like it might
be caused by spaces in the project path. A controlled test disproved it, and the
same test also confirmed that CMake, CUDA and latexmk all work in a path with
spaces, which was worth knowing.

Wrote the dash linter first, because it gates everything. Its first run found an
em dash in the build specification itself.

\section{Phase 1: numerics and serial solvers}

Wrote the core types, the backend interface, the Poisson problem, and validated
early with a smoke test before building further. That smoke test paid for itself
immediately: it confirmed the Jacobi contraction factor matched $\cos(\pi h)$ to
six digits, the Jacobi to Gauss Seidel ratio was exactly 2.0, and the
discretisation error constant was 0.823 against the predicted $\pi^2/12 =
0.8225$. Knowing the foundation was right made everything after it cheaper to
debug.

Then the numerics module, the dense problem, the solver zoo, and the test
harness.

\section{Phase 2: convergence tests}

Wrote the tests that check measurements against closed form theory. Three faults
surfaced here, all in solvers rather than tests: Richardson diverging from a bad
default, Richardson diverging from a bad step estimator, and the conjugate
gradient count refusing to grow with the grid. The last of these turned out to be
the eigenvector problem and was the most interesting bug of the project.

\section{Phase 3: shared memory backends}

OpenMP, pthreads and the jthread pool, then the equivalence suite. The
equivalence suite passed on its first run, which I did not expect; the
determinism design carried it. Found the jthread destructor ordering hazard by
inspection while reviewing the shutdown path.

\section{Phase 4: distributed backends}

MPI and hybrid. Four test cases failed at once at two and four ranks, all from
one cause: the returned solution was left distributed while the tests compared
the whole vector. Adding a single end of solve gather fixed all four. While
reviewing the dense path I found a second, latent fault: block ranges and row
ranges differ, and a gather assuming otherwise would have corrupted vectors
silently.

\section{Phase 5: CUDA}

The link broke in a way that took two wrong guesses to diagnose. Then duplicate
device stubs from kernels in a header. Then a host pointer passed to a device to
device copy. Then the one bit fused multiply add discrepancy, which was the only
one of the four that taught me something about the project rather than about the
tools.

Confirmed the GPU sweeps are bit identical to the CPU, the red black iteration
penalty is 2.6 percent, and the device triad bandwidth is around 548 GiB/s.

\section{Phase 6: the sweep}

Declared the matrix, then found two faults in my own declaration before running
it: two quadratic methods listed in a block meant for linear ones, and the resume
logic testing on the wrong side of the work. Both were caught by thinking about
the arithmetic rather than by waiting for the consequences.

\section{Phases 7 to 9}

Documentation written from the measured numbers, the three reports, and final
quality assurance.

% =====================================================================
\chapter{Problems and errors}

The full record is in \texttt{docs/ENGINEERING\_LOG.md} and in the debug report,
with symptom, root cause, options, fix and verification for each. This chapter
records what I want to remember about them rather than repeating them.

\section{The five that produced plausible numbers}

These are the ones worth internalising. A crash teaches nothing; a wrong answer
that looks right is the failure mode that survives.

\paragraph{The manufactured solution was an eigenvector.} With
$u = \sin(\pi x)\sin(\pi y)$, conjugate gradient converged in one iteration at
every grid size, because the source is exactly the lowest eigenvector of the
stencil and the Krylov subspace is therefore one dimensional. It looked like a
spectacular result.

What caught it was a test asserting the iteration count should \emph{grow} with
the grid, not merely that the solver converged. I would not have written that
test if the specification had not asked for growth orders to be verified.

The fix needed thought, because the obvious alternatives are worse: a polynomial
manufactured solution has vanishing fourth derivatives, which makes the stencil
exact and destroys the second order accuracy test; a sum of a few modes just
moves the problem. Keeping both right hand sides and using each where it is
honest was the right answer, and the single mode source is genuinely \emph{ideal}
for the stationary methods, because the slowest Jacobi mode is that same lowest
mode.

\paragraph{A default that suited the first solver written.}
\texttt{relaxation} defaulted to 1.0, harmless for SOR and catastrophic for
Richardson. The generalisation: a shared options struct wants defaults that mean
"ask the component", not defaults that happen to suit one of them.

\paragraph{The resumable sweep was not resumable.} The skip test sat after the
work. Counters reported plausible numbers throughout. It survived until I timed a
restart instead of reading its output.

\paragraph{A tolerance below the arithmetic.} Two tests asked for $10^{-13}$ from
stationary iterations whose rounding floor on this operator is around
$3\times10^{-13}$. The failure appeared only when the solver those tests used
changed, so it looked like a regression in the solver. Measuring the floor turned
a confusing symptom into a table.

\paragraph{One bit of fused multiply add.} The host was contracting $ab + cd$ and
the device was not. Visible only because the equivalence suite asserts bit
identity. Following it led to a real weakness in the project's central claim, and
the fix, disabling contraction on both sides, made the claim robust to compiler
upgrades as well.

\section{The toolchain problems}

Four of these, all in phase 0 or phase 5, all structural: nvcc rejecting the
project compiler, a CMake setting arriving after \texttt{project()}, MPI version
macros describing conformance rather than capability, and the CUDA host
compiler's library directory hijacking the link.

The last one cost the most time and taught the most: I guessed twice before
diagnosing it, first blaming spaces in the path and then blaming the linker
driver. The lesson is that when a link fails on symbols from a component that has
nothing to do with the change, the change probably altered the \emph{search
order} rather than the code.

\section{The one the environment imposed}

WSL2 hides the hybrid core topology. I could not identify performance versus
efficiency cores from inside the guest, and thread affinity binds to a virtual
processor the hypervisor may place anywhere.

The response I am happiest with in the whole project: rather than assume the
conventional processor enumeration and produce confident labels with nothing
behind them, the library measures, reports what it could establish, and takes the
knee from the aggregate scaling curve, which does not depend on identifying
individual cores. The claim got weaker and became true.

% =====================================================================
\chapter{Future work}

\section{Preconditioned conjugate gradient}

The most obvious gap. Symmetric Gauss Seidel and SSOR are already implemented and
are exactly the preconditioners the method wants; they currently exist in the zoo
without the application that motivates them. This is a small amount of code and
would substantially improve the strongest solver in the library.

\section{Multigrid}

Every component is present: smoothers in the relaxation sweeps, an exact coarse
solve in the LU routine, the grid hierarchy implicit in the Poisson problem.
Multigrid makes the iteration count independent of the grid, which would put the
$O(n^2)$ and $O(n)$ results of this project in their proper context as the two
things it improves on. It deserves its own project.

\section{A local slab MPI decomposition}

Each rank currently allocates the whole grid. The communication volume is already
correct so no timing here would change, but it would remove the memory ceiling
that keeps distributed runs below the sizes the device comparison uses.

\section{Bare metal measurement of the core knee}

The one objective the environment prevented. Running the same sweep outside a
hypervisor would settle whether the knee in the aggregate curve coincides with
the performance core count, which is currently inferred rather than confirmed.

\section{Mixed precision}

The stencil sweeps are bandwidth bound, so halving the width of the iterate would
come close to halving the time. Whether the convergence penalty is worth it is
exactly the kind of question this library's separation of iteration count from
per iteration cost is built to answer, and it would also test whether the
determinism design survives a second scalar type.

\section{Smaller things worth doing}

\begin{itemize}[leftmargin=*]
  \item \textbf{An allocation free block Jacobi.} It currently snapshots the
        previous iterate each sweep, which is an $O(N)$ copy that could be
        avoided by double buffering.
  \item \textbf{Overlapping halo exchange with interior computation.} The
        standard optimisation, and the library is structured to make it
        expressible: compute the interior while the halo is in flight, then the
        boundary rows. It would improve MPI scaling and the measurement
        infrastructure to prove it already exists.
  \item \textbf{Persistent MPI collectives.} The probe confirmed
        \texttt{MPI\_Allreduce\_init} links even though the version macro says
        3.1. Worth measuring whether it helps the conjugate gradient reduction
        cost.
  \item \textbf{A device resident multi solve mode.} Transfer time is currently
        reported separately and honestly, but a mode that keeps state resident
        across solves would show what a real application would see.
\end{itemize}

\end{document}
